\documentclass[11pt]{amsart}

\usepackage[T1]{fontenc}
\usepackage{lmodern}
\usepackage{microtype}
\usepackage{mathtools,amssymb,amsthm}
\usepackage[hidelinks]{hyperref}

\hypersetup{
  pdftitle={K5 minus an edge is 2-conformally rigid but not 1-conformally rigid}
}

\newtheorem{theorem}{Theorem}
\newtheorem{lemma}[theorem]{Lemma}
\newtheorem{proposition}[theorem]{Proposition}
\newtheorem{corollary}[theorem]{Corollary}
\theoremstyle{remark}
\newtheorem*{remark}{Remark}

\DeclareMathOperator{\tr}{tr}
\newcommand{\Real}{\mathbb R}
\newcommand{\DeltaE}{\Delta_E}
\newcommand{\Lop}{L}
\newcommand{\Ladj}{L^{*}}

\title[$K_5-e$ is $2$- but not $1$-conformally rigid]
{$K_5$ Minus an Edge is $2$-Conformally Rigid but not $1$-Conformally Rigid}

\date{}

\subjclass[2020]{05C50, 05C22, 90C22, 15A42}
\keywords{conformal rigidity, weighted Laplacian, eigenvalue sums, Ky Fan principle,
dual certificate, counterexample}

\begin{document}

\begin{abstract}
Let $G$ be a graph on $n$ vertices with edge set $E$, let
$\DeltaE=\{w\in\Real_{\ge0}^{E}:\sum_e w_e=|E|\}$, and index the eigenvalues of the
weighted Laplacian $\Lop(w)$ increasingly as
$0=\lambda_1(w)\le\cdots\le\lambda_n(w)$. Write $S_k(w)$ for the sum of the $k$ largest
and $s_k(w)$ for the sum of the $k$ smallest nontrivial eigenvalues. Following
Assump\c{c}\~ao, Coutinho and Godsil, $G$ is \emph{$k$-conformally rigid} if the all-ones
weight simultaneously minimises $S_k$ and maximises $s_k$ on $\DeltaE$. In the Conclusion
of their paper they ask whether there is a graph that is $2$-conformally rigid but not
$1$-conformally rigid. There is one at $n=5$: $K_5$ minus an edge. We answer the question with an
exact rational dual certificate, a single $5\times5$ integer matrix, that proves
$S_2(w)\ge10=S_2(\mathbf 1)$ for every $w\in\DeltaE$; since $n-1-k=2=k$ here, $2$-conformal
rigidity follows. At this order the duality is self-paired, so the two conditions defining
$2$-conformal rigidity coincide and one inequality discharges both; a witness at larger $n$,
where the two conditions are genuinely distinct, is not provided here. That $K_5-e$ is not $1$-conformally rigid is already recorded as
numerical data by Niu; we exhibit the exact rational weight behind it.
\end{abstract}

\maketitle

\section{The question, and the answer}

Let $G=(V,E)$ be a graph on $n$ vertices. For $w\in\Real_{\ge0}^{E}$ let
\[
 \Lop(w)=\sum_{\{a,b\}\in E}w_{ab}\,(e_a-e_b)(e_a-e_b)^{\mathsf T}
\]
be the weighted Laplacian, and let its eigenvalues be indexed \emph{increasingly},
$0=\lambda_1(w)\le\lambda_2(w)\le\cdots\le\lambda_n(w)$. Put
\[
 \DeltaE=\Bigl\{w\in\Real_{\ge0}^{E}\ :\ \textstyle\sum_{e\in E}w_e=|E|\Bigr\},
\]
\[
 S_k(w)=\sum_{j=n-k+1}^{n}\lambda_j(w),
 \qquad
 s_k(w)=\sum_{j=2}^{k+1}\lambda_j(w),
\]
so that $S_k$ and $s_k$ sum the $k$ largest and the $k$ smallest \emph{nontrivial}
eigenvalues. The all-ones weight $\mathbf 1$ lies in $\DeltaE$ and $\Lop(\mathbf 1)$ is the
ordinary Laplacian. Following the definition of $k$-conformal rigidity in
Assump\c{c}\~ao, Coutinho and Godsil \cite{ACG}, call $G$
\emph{upper-$k$ rigid} if $\mathbf 1$ minimises $S_k$
over $\DeltaE$, \emph{lower-$k$ rigid} if $\mathbf 1$ maximises $s_k$ over $\DeltaE$, and
\emph{$k$-conformally rigid} if both hold, for $k\in\{1,\dots,n-1\}$. For $k=1$ this is the
notion of Steinerberger and Thomas \cite{ST}: $\mathbf 1$ minimises $\lambda_n$ and
maximises $\lambda_2$; see also Gouveia, Steinerberger and Thomas \cite{GST} for the
$k=1$ theory.

The Conclusion of \cite{ACG} asks, in running prose,

\begin{quote}
``Are there graphs that are $2$-conformally rigid but not $1$-conformally rigid?''
\end{quote}

\noindent
The answer is yes, at $n=5$.

\begin{theorem}\label{thm:main}
Let $H=K_5-e$ be the complete graph on $V=\{1,2,3,4,5\}$ with the pair $\{4,5\}$ deleted,
so $|E(H)|=9$. Then $H$ is $2$-conformally rigid and is not $1$-conformally rigid.
\end{theorem}

Both halves are certified by objects printed below in full.

The negative half of Theorem~\ref{thm:main} is not new as a fact: Niu's dataset
\cite{Niu} tabulates lower- and upper-$1$ rigidity for the graphs of the House of Graphs
census, and its row \texttt{id=450}, the $5$-vertex $9$-edge graph $H$, records
\texttt{lcr=False}, \texttt{ucr=True}. What is supplied here is the dual certificate of
Proposition~\ref{prop:upper2}, and the exact rational weight $w^{*}$ of
Proposition~\ref{prop:lower1} behind that table cell. No claim of priority is made for
the upper-$2$ fact itself: whether it already follows from results or data in \cite{ACG}
or \cite{Niu} has not been checked here. Note that \cite{Niu} calls a graph
conformally rigid when it is lower-$1$ \emph{or} upper-$1$ rigid, while \cite{ST} and
\cite{ACG} use \emph{and}; we use the conjunction of \cite{ACG} throughout.

\section{The two certificates}

Both proofs use the adjoint of $w\mapsto\Lop(w)$. For a symmetric $X$ put
\begin{equation}\label{eq:adj}
\begin{gathered}
 \Ladj(X)_{ab}=X_{aa}+X_{bb}-2X_{ab}\qquad(\{a,b\}\in E),\\
 \text{so that}\qquad
 \tr\bigl(\Lop(w)X\bigr)=\sum_{e\in E}w_e\,\Ladj(X)_e ,
\end{gathered}
\end{equation}
which is immediate from
$\tr\bigl((e_a-e_b)(e_a-e_b)^{\mathsf T}X\bigr)=(e_a-e_b)^{\mathsf T}X(e_a-e_b)$.

The spectrum of $H$ is exact and elementary. Since $\Lop(\mathbf 1)=5I-J-dd^{\mathsf T}$
with $J$ the all-ones matrix and $d=e_4-e_5$, a vector $x$ with $x^{\mathsf T}\mathbf 1=0$
and $x_4=x_5$ satisfies $\Lop(\mathbf 1)x=5x$, while $\Lop(\mathbf 1)d=3d$. Hence
\begin{equation}\label{eq:spec}
 \operatorname{spec}\Lop(\mathbf 1)=\{0,3,5,5,5\},
 \qquad
 S_1=5,\ \ S_2=10,\ \ S_3=15,\ \ S_4=18,\ \ s_1=3 ,
\end{equation}
the $5$-eigenspace being $W=\{x:x^{\mathsf T}\mathbf 1=0,\ x_4=x_5\}$, spanned by
$a=(1,-1,0,0,0)$, $b=(0,1,-1,0,0)$ and $c=(2,2,2,-3,-3)$.

\begin{proposition}[upper-$2$ rigidity]\label{prop:upper2}
Let
\begin{equation}\label{eq:M}
 X=\frac{1}{45}M,
 \qquad
 M=\begin{pmatrix}
  22 & -3 & -3 & -8 & -8\\
  -3 & 22 & -3 & -8 & -8\\
  -3 & -3 & 22 & -8 & -8\\
  -8 & -8 & -8 & 12 & 12\\
  -8 & -8 & -8 & 12 & 12
 \end{pmatrix}.
\end{equation}
Then $S_2(w)\ge10=S_2(\mathbf 1)$ for every $w\in\DeltaE$, so $H$ is upper-$2$ rigid.
\end{proposition}

\begin{proof}
Four facts about $X$, each a direct computation with \eqref{eq:M}.

\emph{(1) $X$ is symmetric with $\tr X=(3\cdot22+2\cdot12)/45=90/45=2$.}

\emph{(2) $X\mathbf 1=0$ and $Xd=0$, and $Xa=\tfrac59a$, $Xb=\tfrac59b$,
$Xc=\tfrac89c$.} The first two hold because every row of $M$ sums to $0$ and its last two
columns are equal. For the others, $Ma=(25,-25,0,0,0)$, which is $25a$, and
$Mc=(80,80,80,-120,-120)$, which is $40c$; and $Mb=25b$ because swapping the coordinates
$1$ and $3$ permutes $M$ to itself and carries $a$ to $-b$. Since $\mathbf 1,d,a,b,c$ are linearly
independent (the determinant of the matrix with these columns is $-30$), the spectrum of
$X$ is exactly $\{0,0,\tfrac59,\tfrac59,\tfrac89\}$; equivalently the characteristic
polynomial of $X$ is $x^2(x-\tfrac89)(x-\tfrac59)^2$. In particular
\[
 0\preceq X\preceq I
\]
holds \emph{exactly}, not to a tolerance.

\emph{(3) $\Lop(\mathbf 1)X=5X$, hence $\tr(\Lop(\mathbf 1)X)=5\tr X=10=S_2(\mathbf 1)$.}
Indeed the range of $X$ is spanned by $a,b,c$ and lies in the $5$-eigenspace $W$.

\emph{(4) $\Ladj(X)$ is the constant $\tfrac{10}{9}$ on all nine edges.} For an edge inside
$\{1,2,3\}$, $\Ladj(X)=\tfrac{1}{45}(22+22+6)=\tfrac{50}{45}=\tfrac{10}{9}$; for an edge
from $\{1,2,3\}$ to $\{4,5\}$, $\Ladj(X)=\tfrac{1}{45}(22+12+16)=\tfrac{50}{45}=\tfrac{10}{9}$.
These are the only two edge types, so
$\sum_{e\in E}\Ladj(X)_e=9\cdot\tfrac{10}{9}=10$.

Now let $w\in\DeltaE$. Because $X\mathbf 1=0$, $X$ acts on the $(n-1)$-dimensional space
$\mathbf 1^{\perp}$, on which $\Lop(w)$ has eigenvalues $\lambda_2(w),\dots,\lambda_n(w)$.
The Ky Fan maximum principle \cite{Fan} applied on that subspace gives
\[
 \tr\bigl(\Lop(w)X\bigr)\le S_2(w)
 \qquad\text{whenever } 0\preceq X\preceq I,\ \tr X=2,\ X\mathbf 1=0 .
\]
By (1), (2), \eqref{eq:adj} and (4),
\[
 \tr\bigl(\Lop(w)X\bigr)=\sum_{e\in E}w_e\,\Ladj(X)_e=\frac{10}{9}\sum_{e\in E}w_e
 =\frac{10}{9}\,|E|=10 ,
\]
the constancy of $\Ladj(X)$ being exactly what makes this independent of $w$. Hence
$S_2(w)\ge10$ for every $w\in\DeltaE$, and by (3) equality holds at $w=\mathbf 1$.
\end{proof}

\begin{lemma}[duality]\label{lem:dual}
For $1\le k\le n-2$ and every $w\in\DeltaE$ we have $S_k(w)+s_{n-1-k}(w)=2|E|$.
Consequently $G$ is upper-$k$ rigid if and only if it is lower-$(n-1-k)$ rigid.
\end{lemma}

\begin{proof}
$\sum_{j=2}^{n}\lambda_j(w)=\tr\Lop(w)=2\sum_ew_e=2|E|$ on $\DeltaE$, and $S_k$ and
$s_{n-1-k}$ partition that sum. A constant sum turns minimising one summand into
maximising the other.
\end{proof}

This is the duality of \cite{ACG}, stated there for the same range
$k\in\{1,\dots,n-2\}$, which is what $s_{n-1-k}$ requires.

\begin{proposition}[failure of lower-$1$ rigidity]\label{prop:lower1}
Let $w^{*}\in\Real_{\ge0}^{E}$ take the value $\tfrac97$ on the six edges meeting
$\{4,5\}$ and $\tfrac37$ on the three edges inside $\{1,2,3\}$. Then $w^{*}\in\DeltaE$,
\[
 \operatorname{spec}\Lop(w^{*})=\Bigl\{0,\ \tfrac{27}{7},\ \tfrac{27}{7},\ \tfrac{27}{7},\ \tfrac{45}{7}\Bigr\},
\]
and $s_1(w^{*})=\tfrac{27}{7}>3=s_1(\mathbf 1)$. Hence $H$ is not lower-$1$ rigid, so it is
not $1$-conformally rigid; equivalently, $S_3(w^{*})=\tfrac{99}{7}<15=S_3(\mathbf 1)$, so
$H$ is not upper-$3$ rigid either.
\end{proposition}

\begin{proof}
$w^{*}$ is nonnegative and $6\cdot\tfrac97+3\cdot\tfrac37=\tfrac{54+9}{7}=9=|E|$, so
$w^{*}\in\DeltaE$. Write $\alpha=\tfrac37$, $\beta=\tfrac97$. With $a,b,c,d$ as above:
$\Lop(w^{*})a=(3\alpha+2\beta)a=\tfrac{9+18}{7}a=\tfrac{27}{7}a$ and likewise for $b$;
$\Lop(w^{*})d=3\beta d=\tfrac{27}{7}d$, because $d$ is supported on $\{4,5\}$, which is a
non-edge, and each of $4,5$ has weighted degree $3\beta$; and reading off the first and
fourth coordinates of $\Lop(w^{*})c$ gives $\tfrac{90}{7}=\tfrac{45}{7}\cdot2$ and
$-\tfrac{135}{7}=\tfrac{45}{7}\cdot(-3)$, so $\Lop(w^{*})c=\tfrac{45}{7}c$. Together with
$\Lop(w^{*})\mathbf 1=0$ these five independent vectors give the stated spectrum, whose sum
is $\tfrac{81+45}{7}=18=2|E|$ as it must be. Now
$s_1(w^{*})=\lambda_2(w^{*})=\tfrac{27}{7}>3=\lambda_2(\mathbf 1)=s_1(\mathbf 1)$, so
$\mathbf 1$ does not maximise $s_1$ on $\DeltaE$. The upper-$3$ statement is the same fact
through Lemma~\ref{lem:dual} with $n=5$, $k=3$: $S_3(w^{*})=\tfrac{27+27+45}{7}=\tfrac{99}{7}$
and $18-\tfrac{27}{7}=\tfrac{99}{7}<15$.
\end{proof}

\begin{proof}[Proof of Theorem~\ref{thm:main}]
For $n=5$ and $k=2$ we have $n-1-k=2=k$, so Lemma~\ref{lem:dual} turns upper-$2$ rigidity
into lower-$2$ rigidity. Proposition~\ref{prop:upper2} gives the former, hence both, hence
$H$ is $2$-conformally rigid. Proposition~\ref{prop:lower1} gives that $H$ is not
lower-$1$ rigid, hence not $1$-conformally rigid.
\end{proof}

\begin{remark}
At $n=5$ and $k=2$ the duality of Lemma~\ref{lem:dual} is self-paired, so the two
conditions making up $2$-conformal rigidity coincide and one inequality discharges both.
In general they are distinct, and a witness at larger $n$ would have to verify each; this
note says nothing about that case.
\end{remark}

\section{Verification}

Everything above is determined by the labelled graph $H$ and the two printed objects $M$
and $w^{*}$. The nine edges of $H$ in the labelling of this note are
\[
 \{1,2\},\{1,3\},\{1,4\},\{1,5\},\{2,3\},\{2,4\},\{2,5\},\{3,4\},\{3,5\},
\]
and a machine-readable encoding is the graph6 string \texttt{D\textasciicircum\{}, on the
vertex set $\{0,1,2,3,4\}$ with the pair $\{0,1\}$ deleted, carried onto the labelling
above by $0\mapsto4$, $1\mapsto5$, $2\mapsto1$, $3\mapsto2$, $4\mapsto3$.

The proofs are $5\times5$ rational arithmetic and can be checked by hand. A stand-alone
program \texttt{verify.py} accompanies this note; it reads the graph6 string, the edge
list, the matrix $M$ and the weight $w^{*}$ exactly as printed here, and re-derives in
exact integer and \texttt{Fraction} arithmetic, using only the Python standard library and
with no floating-point comparison anywhere: the graph and its edge list under the printed
relabelling; the characteristic polynomial and spectrum \eqref{eq:spec} together with
$S_1,\dots,S_4$ and $s_1$ at $\mathbf 1$; the pairing identity \eqref{eq:adj} on several
rational weights; every claim (1)--(4) about $X$, including the exact containment
$0\preceq X\preceq I$; and that $w^{*}\in\DeltaE$ with the stated spectrum and the two
strict inequalities $\tfrac{27}{7}>3$ and $\tfrac{99}{7}<15$. Its recorded run reports
$68$ checks, all passing. It does not reprove the Ky Fan principle or
Lemma~\ref{lem:dual}, and it certifies nothing about minimality, about other graphs, or
about other orders.

\begin{thebibliography}{9}

\bibitem{ACG}
R.~Assump\c{c}\~ao, G.~Coutinho, and C.~Godsil,
\emph{Total conformal rigidity in graphs},
arXiv:2605.08508v2 (24 Jul 2026).
\href{https://arxiv.org/abs/2605.08508}{arXiv:2605.08508}.
The definition, the duality and the question quoted above are from v2.

\bibitem{Fan}
K.~Fan,
\emph{On a theorem of Weyl concerning eigenvalues of linear transformations I},
Proc. Nat. Acad. Sci. U.S.A. \textbf{35} (1949), 652--655.
\href{https://doi.org/10.1073/pnas.35.11.652}{doi:10.1073/pnas.35.11.652}.

\bibitem{Niu}
A.~Niu,
\emph{Conformal rigidity of graphs: subdifferentials and orbit-isometries},
arXiv:2605.15017v1, 2026.
\href{https://arxiv.org/abs/2605.15017}{arXiv:2605.15017};
dataset at \href{https://github.com/andrewmniu/conformally-rigid-graphs}{github.com/andrewmniu/conformally-rigid-graphs},
file \texttt{hog/numerical/numerical\_certs.csv}, row \texttt{id=450}.

\bibitem{ST}
S.~Steinerberger and R.~Thomas,
\emph{Conformally rigid graphs},
J. Graph Theory \textbf{109} (2025), no.~3, 366--386.
\href{https://doi.org/10.1002/jgt.23229}{doi:10.1002/jgt.23229};
\href{https://arxiv.org/abs/2402.11758}{arXiv:2402.11758}.

\bibitem{GST}
J.~Gouveia, S.~Steinerberger, and R.~Thomas,
\emph{Conformal rigidity and spectral embeddings of graphs},
Ann. Comb. (2026), \href{https://doi.org/10.1007/s00026-026-00828-8}{doi:10.1007/s00026-026-00828-8};
\href{https://arxiv.org/abs/2506.20541}{arXiv:2506.20541}.

\end{thebibliography}

\end{document}
